\chapter{Other Non-Functional Requirements}

\section{Performance Requirements}
\begin{itemize}[leftmargin=*]
    \item \textbf{Browser Startup Overhead:} The CDP switch injection at startup adds negligible overhead (single conditional check on the command-line object) with no measurable impact on browser launch time.
    \item \textbf{AI Command Latency:} Simple single-step commands (e.g., ``open a new tab'') shall execute within 2 seconds on a machine with a modern LLM running locally. Multi-step agentic tasks may take longer depending on the complexity of the plan and network page load times.
    \item \textbf{LLM Inference Speed:} For smooth interaction, token generation by the local LLM shall sustain at least 10 tokens/second on recommended hardware (8 GB RAM, modern CPU).
    \item \textbf{CDP Round-Trip Latency:} CDP commands sent over the loopback WebSocket shall complete with latency under 10 ms, excluding page-load wait time.
    \item \textbf{UI Responsiveness:} The browser UI (toolbar, tabs, address bar) shall remain fully responsive during AI agent processing, as the agent runs in a separate Python process.
\end{itemize}

\section{Safety Requirements}
\begin{itemize}[leftmargin=*]
    \item The AI agent shall not execute browser actions that could cause irreversible data loss (e.g., deleting browser profiles) without explicit user confirmation.
    \item The CDP interface is bound to the loopback address (\texttt{127.0.0.1}) only; it must not be exposed on any network interface accessible from the internet.
    \item AI-generated JavaScript executed via \texttt{Runtime.evaluate} shall be sandboxed within the browser's existing renderer security model.
\end{itemize}

\section{Security Requirements}
\begin{itemize}[leftmargin=*]
    \item \textbf{No Data Exfiltration:} The AI agent shall establish no connections to external servers for AI processing. All LLM inference and STT processing shall be performed on-device.
    \item \textbf{Loopback Binding:} Both the AI agent server (port 7788) and the CDP endpoint (port 9222) shall accept connections from \texttt{127.0.0.1} only.
    \item \textbf{Session Isolation:} The AI agent shall not access content from other browser profiles or user sessions without explicit permission.
    \item \textbf{Input Sanitisation:} User commands passed to the LLM shall be sanitised to prevent prompt injection attacks that could cause the model to generate harmful CDP instructions.
    \item \textbf{Chromium Security Baseline:} All security features inherited from Chromium (sandboxed renderer processes, site isolation, safe browsing) shall remain fully active in Prism.
\end{itemize}

\section{Software Quality Attributes}

\subsection{Reliability}
\begin{itemize}[leftmargin=*]
    \item The browser shall be at least as stable as the underlying Chromium base revision on which it is built.
    \item If the AI agent server is not running, clicking the Prism AI button shall open a new tab showing a connection error page; the rest of the browser shall function normally.
    \item If a CDP command fails or times out, the AI agent shall report an informative error to the UI and allow the user to retry.
\end{itemize}

\subsection{Maintainability}
\begin{itemize}[leftmargin=*]
    \item All Prism-specific code shall be clearly separated from upstream Chromium code through focused, well-described patch files.
    \item Patch files shall be regenerated against a new Chromium base revision when the upstream drifts significantly, using the procedure documented in the \texttt{README.md}.
    \item New features shall each be captured in separate, logically named patch files (e.g., \texttt{0003-voice-input.diff}) to enable independent review and rollback.
\end{itemize}

\subsection{Usability}
\begin{itemize}[leftmargin=*]
    \item The Prism AI Button shall be immediately identifiable in the toolbar with a distinct icon and a clear ``Prism AI'' tooltip on hover.
    \item The AI agent chat interface shall provide clear feedback: a progress indicator during multi-step task execution, and colour-coded success/failure messages on completion.
    \item Voice commands shall be acknowledged with a visual recording indicator in the UI.
    \item Users with no technical knowledge shall be able to perform basic AI-assisted browsing tasks without consulting documentation.
\end{itemize}

\subsection{Portability}
\begin{itemize}[leftmargin=*]
    \item Prism shall build and run on all platforms supported by Chromium: Linux (primary), macOS, and Windows 10 or later.
    \item The patch set shall be platform-independent where possible; platform-specific code sections in the patches shall use Chromium's \texttt{BUILDFLAG(IS\_WIN)}, \texttt{BUILDFLAG(IS\_MAC)}, and \texttt{BUILDFLAG(IS\_LINUX)} guards.
\end{itemize}

\subsection{Extensibility}
\begin{itemize}[leftmargin=*]
    \item The AI agent backend shall be designed with a modular tool interface, allowing new browser automation capabilities (e.g., screenshot capture, PDF export) to be added as separate tool modules without modifying core agent logic.
    \item Chromium extensions that interact with the Prism AI agent via the standard Chrome Extension Messaging API shall be supportable.
\end{itemize}

\section{Business Rules}
\begin{itemize}[leftmargin=*]
    \item Prism AI Browser is developed for educational and research purposes as a mini project.
    \item The project is open-source and its patch repository is publicly available on GitHub.
    \item The browser may not be distributed commercially without compliance with the Chromium / BSD licence terms.
\end{itemize}
